\subsection{Modular continuations and observation schedules}
\label{sec:strategies}
The \texttt{inline} strategy above tests states reached by executing
available bodies and loop iterations. The \texttt{modular} strategy instead
uses contracts, invariants, and assertions to construct continuations.
At a boundary it checks the incoming obligation before abstracting and
sampling a legal continuation; sampling does not erase an earlier obligation.
Both strategies reuse QCP's resource matching, invocation bindings and frames.

For a modular call, the contract constrains the returned value and modified
resources. A modular loop checks entry, samples a guard-satisfying invariant
state for a preservation path, and samples an invariant state satisfying the
exit guard for the continuation. These need not be successive states of an
actual loop run. An assertion can similarly provide a boundary state.
Sampling retains the represented relational constraints and framed resources;
it cannot reconstruct constraints omitted from the annotation.
Unsupported or unsuccessful sampling remains unresolved.

The observation schedule is independent of this choice.
\texttt{entry-exit} selects entry, safety and final checks;
\texttt{all-boundaries} additionally selects internal observations.
A call without a body uses contract instantiation in either strategy.
A declaration-only root provides an entry check, not evidence of successful
execution or a satisfied postcondition.

\paragraph{Interpreting the feedback.}
Body execution checks reached states, whereas modular execution checks
selected models of summaries. For example, let a helper return 1 under the
contract $0\le\mathit{ret}\le1$, and let its caller require $\mathit{ret}=1$.
Inline execution passes this check; modular execution may sample 0 and
expose that the contract does not suffice for the caller.
This diagnoses an interface obligation without exhibiting a violating
execution of the available helper body. Archived paired controls exercise
both this case and an exact-summary case under both observation schedules.

Reports therefore identify the transition model and the states covered.
A passing sample does not establish all behaviors admitted by a summary.
Mixed executions with bodyless calls inherit this qualification unless
those external transitions have a justified implementation model.
Section~\ref{sec:mode-correspondence} relates the two strategies when the
complete modular specification has already been verified.
